\subsubsection{Comparison across Pauli noise channels}

{ In Tables~I--V, boldface marks the lowest susceptibility in each column, while
italics mark the second-lowest value; ties can occur after rounding.


  } \begin{table}[tbp]
    \centering
    \caption{Noise susceptibility $\chi$ for H2, 4 qubits.}
    \begin{adjustbox}{max width=\linewidth}
    \begin{tabular*}{\linewidth}{@{\extracolsep{\fill}} l c c c c c @{}}
    \toprule
    \textbf{Method} & \textbf{D} & \textbf{X} & \textbf{Y} & \textbf{Z} & \textbf{\# CX gates} \\
    \midrule
    JW & $49.1$ & $21.6$ & $40.2$ & $46.7$ & $49$ \\
    BK & $38.1$ & $34.3$ & $30.9$ & $31.0$ & $35$ \\
    JW GdBM & $24.2$ & $\mathit{16.4}$ & $23.2$ & $19.0$ & $21$ \\
    BK GdBM & $19.1$ & $19.6$ & $19.7$ & $\mathbf{0.7}$ & $\mathbf{15}$ \\
    FSN a-t-a & $\mathbf{14.9}$ & $\mathit{17.2}$ & $15.1$ & $\mathit{7.9}$ & $15$ \\
    FSN $2\times N$ & $19.5$ & $21.1$ & $18.9$ & $10.8$ & $21$ \\
    \midrule
    MSN $2\times N$ & $\mathit{18.3}$ & $\mathbf{13.9}$ & $\mathbf{7.3}$ & $13.6$ & $16$ \\
    \bottomrule
    \end{tabular*}
    \end{adjustbox}
    \end{table}


    \begin{table}[tbp]
    \centering
    \caption{Noise susceptibility $\chi$ for H2, 8 qubits.}
    \begin{adjustbox}{max width=\linewidth}
    \begin{tabular*}{\linewidth}{@{\extracolsep{\fill}} l c c c c c @{}}
    \toprule 
    \textbf{Method} & \textbf{D} & \textbf{X} & \textbf{Y} & \textbf{Z} & \textbf{\# CX gates} \\
    \midrule
    JW & $715$ & $297$ & $612$ & $555$ & $417$ \\
    BK & $902$ & $567$ & $796$ & $735$ & $394$ \\
    JW GdBM & $487$ & $394$ & $452$ & $400$ & $151$ \\
    BK GdBM & $460$ & $369$ & $385$ & $278$ & $\mathit{132}$ \\
    FSN a-t-a & $217$ & $\mathit{255}$ & $\mathit{218}$ & $\mathbf{109}$ & $170$ \\
    FSN $2\times N$ & $\mathit{267}$ & $305$ & $254$ & $135$ & $\mathbf{96}$ \\
    \midrule
    MSN $2\times N$ & $\mathbf{171}$ & $\mathbf{126}$ & $\mathbf{103}$ & $\mathit{124}$ & $134$ \\
    \bottomrule
    \end{tabular*}
    \end{adjustbox}
    \end{table}
    \begin{table}[tbp]
    \centering
    \caption{Noise susceptibility $\chi$ for LiH, 8 qubits.}
    \begin{adjustbox}{max width=\linewidth}
    \begin{tabular*}{\linewidth}{@{\extracolsep{\fill}} l c c c c c @{}}
    \toprule 
    \textbf{Method} & \textbf{D} & \textbf{X} & \textbf{Y} & \textbf{Z} & \textbf{\# CX gates} \\
    \midrule
    JW & $634$ & $220$ & $531$ & $616$ & $417$ \\
    BK & $696$ & $409$ & $591$ & $610$ & $394$ \\
    JW GdBM & $539$ & $386$ & $477$ & $464$ & $151$ \\
    BK GdBM & $494$ & $369$ & $427$ & $309$ & $\mathit{132}$ \\
    FSN a-t-a & $\mathit{230}$ & $\mathit{293}$ & $\mathit{230}$ & $\mathbf{96}$ & $170$ \\
    FSN $2\times N$ & $282$ & $346$ & $269$ & $\mathit{125}$ & $134$ \\
    \midrule
    MSN $2\times N$ & $\mathbf{205}$ & $\mathbf{149}$ & $\mathbf{123}$ & $146$ & $\mathbf{96}$ \\
    \bottomrule
    \end{tabular*}
    \end{adjustbox}
    \end{table}

    \begin{table}[tbp]
    \centering
    \caption{Noise susceptibility $\chi$ for LiH, 10 qubits.}
    \begin{adjustbox}{max width=\linewidth}
    \begin{tabular*}{\linewidth}{@{\extracolsep{\fill}} l c c c c c @{}}
    \toprule 
    \textbf{Method} & \textbf{D} & \textbf{X} & \textbf{Y} & \textbf{Z} & \textbf{\# CX gates} \\
    \midrule
    JW & $1022$ & $\mathit{344}$ & $866$ & $917$ & $749$ \\
    BK & $1484$ & $825$ & $1269$ & $1184$ & $898$ \\
    JW GdBM & $1632$ & $1258$ & $1377$ & $1368$ & $434$ \\
    BK GdBM & $1773$ & $1462$ & $1578$ & $1162$ & $463$ \\
    FSN a-t-a & $\mathit{318}$ & $408$ & $\mathit{334}$ & $\mathbf{114}$ & $\mathit{216}$ \\
    FSN $2\times N$ & $393$ & $484$ & $387$ & $\mathit{155}$ & $276$ \\
    \midrule
    MSN $2\times N$ & $\mathbf{293}$ & $\mathbf{209}$ & $\mathbf{142}$ & $206$ & $\mathbf{160}$ \\
    \bottomrule
    \end{tabular*}
    \end{adjustbox}
    \end{table}


    \begin{table}[tbp]
    \centering
    \caption{Noise susceptibility $\chi$ for LiH, 12 qubits.}
    \begin{adjustbox}{max width=\linewidth}
    \begin{tabular*}{\linewidth}{@{\extracolsep{\fill}} l c c c c c @{}}
    \toprule 
    \textbf{Method} & \textbf{D} & \textbf{X} & \textbf{Y} & \textbf{Z} & \textbf{\# CX gates} \\
    \midrule
    JW & -- & -- & -- & -- & -- \\
    BK & -- & -- & -- & -- & -- \\
    JW GdBM & $1504$ & $1231$ & $1265$ & -- & $403$ \\
    BK GdBM & $1669$ & $1206$ & -- & -- & $425$ \\
    FSN a-t-a & $\mathit{426}$ & $\mathit{554}$ & $\mathit{429}$ & $\mathbf{159}$ & $321$ \\
    FSN $2\times N$ & $525$ & $653$ & $502$ & $\mathit{212}$ & $411$ \\
    \midrule
    MSN $2\times N$ & $\mathbf{393}$ & $\mathbf{276}$ & $\mathbf{192}$ & $273$ & $\mathbf{240}$ \\
    \bottomrule
    \end{tabular*}
    \end{adjustbox}
    \end{table}


The results in Tables~I--V show that MSN circuits on $2\times N$ architectures generally exhibit
the lowest overall noise susceptibility, although they are not optimal across all error
channels.

Notice that for the $4$-qubit ansatz (H$_2$ molecule), performance is divided among methods: MSN
circuits show the strongest suppression of $X$- and $Y$-type noise, FSN circuits perform best
under $Z$-noise, while BK GdBM unexpectedly exhibits very low susceptibility to $Z$ errors.
However, as the number of qubits increases, the advantage of MSN becomes more pronounced: they
consistently achieve the lowest $\chi$ values for $X$ and $Y$ errors. Interestingly, ladder JW
circuits, although generally the noisiest, display comparatively reduced susceptibility to
$X$-type noise --- sometimes approaching the performance of significantly smaller MSN circuits.
At the same time, FSN retains an advantage for $Z$-noise.

These observations suggest that different circuit constructions naturally specialize in
mitigating different types of errors. This can be explained by the stabilizing action of certain
noise processes on specific decompositions. For instance, in the $4$-qubit BK GdBM circuit, $Z$
errors effectively stabilize the reference state, i.e. with zero excitation angles. Therefore,
for small excitations, such errors induce only minor changes in the purity of the state. A
similar stabilizing effect of $Z$ errors on the reference state can also be observed for ladder
decompositions in JW mapping.

\subsubsection{Phenomenological model for superconducting and ion devices}


\begin{table}[tbp]
    \centering
    \caption{\label{tab:super}Noise susceptibility for molecules in superconducting noise model.}
    \begin{adjustbox}{max width=\linewidth}
    \begin{tabular*}{\linewidth}{@{\extracolsep{\fill}} l c c c c c @{}}
    \toprule 
    \textbf{Method} & H$_2$, 4q & H$_2$, 8q & LiH, 8q & LiH, 10q & LiH, 12q\\
    \midrule
    JW & $0.37$ & $6.40$ & $5.33$ & $9.09$ & -- \\
    BK & $0.32$ & $7.76$ & $5.82$ & $13.28$ & -- \\
    JW GdBM & $0.20$ & $4.64$ & $4.00$ & $7.62$ & -- \\
    BK GdBM & $\mathbf{0.15}$ & $4.41$ & $4.08$ & $8.74$ & $9.91$ \\
    FSN $2\times N$ & $0.16$ & $2.20$ & $2.33$ & $3.24$ & $2.86$ \\
    \midrule
    MSN $2\times N$ & $0.17$ & $\mathbf{1.69}$ & $\mathbf{1.88}$ & $\mathbf{2.70}$ & $\mathbf{2.40}$ \\
    \bottomrule
    \end{tabular*}
    \end{adjustbox}
\end{table}


\begin{table}[tbp]
    \centering
    \caption{\label{tab:ion}Noise susceptibility for {the} molecules in trapped ion
    noise model.}
    \begin{adjustbox}{max width=\linewidth}
    \begin{tabular*}{\linewidth}{@{\extracolsep{\fill}} l c c c c c @{}}
    \toprule 
    \textbf{Method} & H$_2$, 4q & H$_2$, 8q & LiH, 8q & LiH, 10q & LiH, 12q \\
    \midrule
    JW & $0.36$ & $6.55$ & $5.47$ & $9.52$ & -- \\
    BK & $0.30$ & $8.02$ & $5.97$ & $14.06$ & -- \\
    JW GdBM & $0.19$ & $4.74$ & $3.91$ & $7.71$ & $7.79$ \\
    BK GdBM & $0.14$ & $4.33$ & $4.00$ & $8.83$ & $8.70$ \\
    FSN a-t-a & $\mathbf{0.12}$ & $1.72$ & $1.81$ & $2.49$ & $2.62$ \\
    % FSN $2\times N$ & $0.15$ & $2.08$ & $2.18$ & $3.01$ \\
\midrule
    MSN a-t-a & $0.15$ & $\mathbf{1.51}$ & $\mathbf{1.68}$ & $\mathbf{2.41}$ & $\mathbf{2.19}$ \\
    \bottomrule
    \end{tabular*}
    \end{adjustbox}
\end{table}
For noise analysis we use the noise susceptibility with respect to the parameter $\lambda$:
\begin{equation}
    \chi = \left.\frac{d E}{d\lambda}\right|_{\lambda = 0},
\end{equation}

The tables \ref{tab:super} and \ref{tab:ion} show a clear and consistent ranking: MSN/FSN give
the lowest noise susceptibility across most instances, with GdBM variants outperforming plain
JW/BK. Superconducting and ion models agree on the qualitative ordering, though absolute values
differ due to platform-specific error budgets and connectivity; notably, the ion-optimized FSN
variant (``FSN a-t-a'') attains the best value for H$_2$ (4q) thanks to effective all-to-all
coupling.


\subsection{Error sources in VQE}

We quantify the sensitivity of the VQE energy to small noise amplitudes via the susceptibility
\(\chi\).
Results are reported for two phenomenological hardware models—superconducting (left parts, sc)
and trapped-ion (right parts, ion)—and for three noise partitions:
\textbf{Full} (complete model),
\textbf{1q}$\&$\textbf{2q} (single- and two-qubit depolarization only; transverse relaxation \(T_2\) omitted),
and \textbf{2q} (two-qubit depolarization only). Boldface marks the minimum \(\chi\) in each column.

\begin{table}[tbp]
\centering
\caption{Noise susceptibility $\chi$ for H2, 8 qubits.}
\begin{adjustbox}{max width=\linewidth}
% \begin{tabular}{>{\centering\arraybackslash}p{5em}
%                 >{\centering\arraybackslash}p{4em}
%                 >{\centering\arraybackslash}p{4em}
%                 >{\centering\arraybackslash}p{4em}|
%                 >{\centering\arraybackslash}p{4em}
%                 >{\centering\arraybackslash}p{4em}
%                 >{\centering\arraybackslash}p{4em}}
\begin{tabular*}{\linewidth}{@{\extracolsep{\fill}} l c c c | c c c@{}}
\toprule 
\textbf{Method} & \textbf{sc Full} & \textbf{1q}$\&$\textbf{2q} & \textbf{2q} & \textbf{ion Full} & \textbf{1q}$\&$\textbf{2q} & \textbf{2q} \\
\midrule
JW & $6.40$ & $5.65$ & $4.41$ & $6.60$ & $5.52$ & $4.95$ \\
BK & $7.76$ & $6.70$ & $5.25$ & $8.11$ & $6.56$ & $5.90$ \\
JW GdBM & $4.64$ & $4.05$ & $3.01$ & $5.08$ & $3.85$ & $3.38$ \\
BK GdBM & $4.41$ & $3.90$ & $2.94$ & $4.61$ & $3.73$ & $3.30$ \\
FSN  & $2.20$ & $2.09$ & $1.50$ & $1.72$ & $1.61$ & $1.37$ \\
\midrule
MSN  & $\mathbf{1.69}$ & $\mathbf{1.59}$ & $\mathbf{1.02}$ & $\mathbf{1.69}$ & $\mathbf{1.40}$ & $\mathbf{1.14}$ \\
\bottomrule
\end{tabular*}
\end{adjustbox}
\end{table}

\begin{table}[tbp]
\centering
\caption{Noise susceptibility $\chi$ for LiH, 8 qubits.}
\begin{adjustbox}{max width=\linewidth}
\begin{tabular*}{\linewidth}{@{\extracolsep{\fill}} l c c c | c c c@{}}
\toprule 
\textbf{Method} & \textbf{sc Full} & \textbf{1q}$\&$\textbf{2q} & \textbf{2q} & \textbf{ion Full} & \textbf{1q}$\&$\textbf{2q} & \textbf{2q} \\
\midrule
JW & $5.33$ & $4.60$ & $3.53$ & $5.53$ & $4.45$ & $3.97$ \\
BK & $5.82$ & $5.00$ & $3.86$ & $6.05$ & $4.85$ & $4.34$ \\
JW GdBM & $4.00$ & $3.60$ & $2.61$ & $4.15$ & $3.38$ & $2.93$ \\
BK GdBM & $4.08$ & $3.71$ & $2.73$ & $4.19$ & $3.51$ & $3.07$ \\
FSN & $2.33$ & $2.19$ & $1.57$ & $1.81$ & $1.69$ & $1.43$ \\
\midrule
MSN  & $\mathbf{1.88}$ & $\mathbf{1.77}$ & $\mathbf{1.13}$ & $\mathbf{1.84}$ & $\mathbf{1.56}$ & $\mathbf{1.27}$ \\
\bottomrule
\end{tabular*}
\end{adjustbox}
\end{table}

\begin{table}[tbp]
\centering
\caption{Noise susceptibility $\chi$ for LiH, 10 qubits.}
\begin{adjustbox}{max width=\linewidth}
\begin{tabular*}{\linewidth}{@{\extracolsep{\fill}} l c c c | c c c@{}}
\toprule 
\textbf{Method} & \textbf{sc Full} & \textbf{1q}$\&$\textbf{2q} & \textbf{2q} & \textbf{ion Full} & \textbf{1q}$\&$\textbf{2q} & \textbf{2q} \\
\midrule
JW & $9.09$ & $7.62$ & $5.92$ & $9.52$ & $7.42$ & $6.65$ \\
BK & $13.28$ & $11.00$ & $8.76$ & $14.06$ & $10.85$ & $9.84$ \\
JW GdBM & $7.62$ & $6.62$ & $4.96$ & $7.71$ & $6.32$ & $5.57$ \\
BK GdBM & $8.74$ & $7.53$ & $5.83$ & $8.83$ & $7.32$ & $6.55$ \\
FSN & $3.24$ & $3.06$ & $2.18$  & $2.49$ & $2.34$ & $1.99$  \\
\midrule
MSN  & $\mathbf{2.70}$ & $\mathbf{2.55}$ & $\mathbf{1.63}$ & $\mathbf{2.41}$ & $\mathbf{2.24}$ & $\mathbf{1.83}$ \\
\bottomrule
\end{tabular*}
\end{adjustbox}
\end{table}
